\documentclass[12pt, a4paper]{report}

% ==========================================
% PACKAGES
% ==========================================
\usepackage[utf8]{inputenc} % Encoding
\usepackage[T1]{fontenc}    % Font encoding
\usepackage[italian]{babel} % Italian language support for chapters, TOC, etc.
\usepackage{geometry}       % Margin configuration
\geometry{a4paper, top=3cm, bottom=3cm, left=3cm, right=3cm} % Thesis-style margins
\usepackage{hyperref}       % Clickable links in the Table of Contents
\usepackage{graphicx}       % For inserting images and diagrams
\graphicspath{{images/}}    % configuring the graphicx package & images folder
\usepackage{xcolor}         % For custom colors in code blocksghdg
\usepackage{listings}       % The best package for source code formatting
\usepackage{textcomp}       % For common symbols not in utf8
\usepackage[utf8]{inputenc}
\usepackage{pmboxdraw}

% ==========================================
% CODE BLOCK CONFIGURATION (listings)
% ==========================================
% Define some custom colors for code syntax highlighting
\definecolor{codegreen}{rgb}{0,0.6,0}
\definecolor{codegray}{rgb}{0.5,0.5,0.5}
\definecolor{codepurple}{rgb}{0.58,0,0.82}
\definecolor{backcolour}{rgb}{0.95,0.95,0.92}

% #region foldable
\lstdefinestyle{mystyle}{
    backgroundcolor=\color{backcolour},   
    commentstyle=\color{codegreen},
    keywordstyle=\color{magenta},
    numberstyle=\tiny\color{codegray},
    stringstyle=\color{codepurple},
    basicstyle=\ttfamily\footnotesize,
    breakatwhitespace=false,         
    breaklines=true,                 
    captionpos=b,                    
    keepspaces=true,                 
    numbers=left,                    
    numbersep=5pt,                  
    showspaces=false,                
    showstringspaces=false,
    showtabs=false,                  
    tabsize=2
}
% #endregion

% Define the symbols aliases 
\lstset{style=mystyle,
    literate={°}{{\textdegree}}1
             {à}{{\`a}}1
             {è}{{\`e}}1
             {é}{{\'e}}1
             {ì}{{\`i}}1
             {ò}{{\`o}}1
             {ù}{{\`u}}1
}
%configure the links pakage
\hypersetup{
    %colorlinks=true,
    %linkcolor=blue, % Sets color of internal links
    %filecolor=magenta,      
    %urlcolor=cyan,
    %pdfborder={0 0 0} % remove color around links
}
\newcommand{\inlinecode}[1]{\colorbox{backcolour}{\texttt{\detokenize{#1}}}}%inlineCode formatting


% ==========================================
% DOCUMENT META-DATA
% ==========================================
\title{\textbf{Integrazione Zabbix e Hikvision}\\ \large Documentazione Tecnica e Architetturale}
\author{Brando Calderara\thanks{UNIBS, corso ITID, tirocinio in Secoval}}
\date{\today}

% ==========================================
% BEGIN DOCUMENT
% ==========================================
\begin{document}

% Creates the title page
\maketitle

% Creates the Table of Contents automatically
\tableofcontents
\newpage

\chapter{Introduzione}

Il presente documento illustra il lavoro di progettazione, 
sviluppo e messa in produzione di un nuovo sistema di monitoraggio, 
realizzato durante il periodo di tirocinio di Brando Calderara, per una rete distribuita sul territorio della
Comunità Montana e non solo (29 Comuni). 
Questo sistema si aggancia all'esistente progetto "\emph{SmartCity - Edison}".

Il progetto ha avuto come obiettivo iniziale la migrazione da un precedente sistema di 
monitoraggio (PRTG) 
verso una piattaforma centralizzata, scalabile e \textit{open-source} 
basata su \textbf{Zabbix 7.0 LTS} (Long Term Support). 
Su questa piattaforma si sono poi sviluppate anche altre funzionalità, tra cui: catalogazione,
rilevamento di problemi, acquisizione di dati e invio di messaggi verso altri sistemi.

L'infrastruttura gestita è eterogenea e comprende oltre 1000 dispositivi attivi, 
tra cui telecamere di videosorveglianza (Hikvision, Vigilate), apparati di rete \textit{core} ed 
\textit{edge} (Switch, Access Point Wi-Fi pubblici) e router LTE per raggiungere zone più isolate.

Questa documentazione funge sia da relazione sulle scelte implementative 
(illustrando le scelte architetturali e snippet di codice sorgente) 
sia da manuale, con lo scopo di garantire la totale manutenibilità del sistema 
da parte degli amministratori di rete futuri.

\textbf{Warning: le password e i token dovrebbero essere gestiti tramite variabili d'ambiente (file .env)
 o vault sicuri, e non "hardcodati" negli script o passati in chiaro nei log. Per ora sono hardcodati}

\section{Terminologia}
Per facilitare la lettura e la comprensione del presente manuale anche a personale non strettamente sistemistico, di seguito vengono definiti i termini tecnici principali, i protocolli e i concetti architetturali utilizzati nel progetto.

\subsection{Concetti base di Zabbix}
{\small%makes the text smaller
% Unlike itemize (bullets) or enumerate (numbers), 
%the description environment allows for custom, text-based labels for each 
\begin{description} 
    \item[Zabbix:] è una piattaforma software open-source per il monitoraggio proattivo
     di infrastrutture IT, reti, server, macchine virtuali e servizi cloud. 
     Consente di raccogliere, visualizzare e analizzare metriche in tempo reale 
     (come utilizzo CPU, RAM, rete e spazio disco), inviando avvisi immediati in caso di anomalie.
    \item[Host:] qualsiasi dispositivo fisico o virtuale monitorato dal sistema. 
    Nel nostro caso, una telecamera, uno switch, un Access Point o un router.
    \item[Item:] la singola metrica o il singolo dato raccolto da un Host. 
    Esempi di item sono: la risposta di un comando ping a un host o il tempo di risposta, 
    la risposta ricevuta da un comando, il valore di un sensore fisico o virtuale.
    \item[Trigger:] una regola logica associata a un Item. 
    Es. la telecamera non risponde al ping per 3 minuti, il Trigger "scatta" e genera un allarme (Problem).
    \item[Action:] l'azione automatica che il sistema intraprende quando qualche evento soddisfa alcune condizioni.
    Può essere relativa ai trigger o alla scoperta di un nuovo host.
    \item[Template:] un modello preconfigurato che contiene un insieme logico di 
    Item, Trigger e regole. Viene applicato a più dispositivi per evitare di configurarli manualmente.
    \item[Macro:] variabili integrate in Zabbix (scritte tra parentesi graffe,
    es. \texttt{\{HOST.NAME\}}) che permettono di inserire dinamicamente il nome, 
    l'IP o altri dati all'interno di svariati campi, per esempio le mail di allarme.

\end{description}

\subsection{Reti, protocolli e architettura}
\begin{description}
    %\item[CMDB (Configuration Management Database):] Il database che contiene le informazioni autoritative sui componenti informatici di un'organizzazione. Nel nostro progetto, il ruolo di CMDB è ricoperto dai file Excel comunali.
    %\item[Source of Truth (Sorgente della Verità):] Un concetto architetturale per cui esiste un unico punto di riferimento (il CMDB) considerato come il dato ufficiale e corretto, a cui tutti gli altri sistemi (incluso Zabbix) devono allinearsi.
    \item[SNMP (Simple Network Management Protocol):] 
    il linguaggio standard universale utilizzato dai dispositivi di rete per comunicare
    il proprio stato di salute e la propria configurazione a un server di monitoraggio centrale.
    \item[MIB e OID:] Il \textbf{MIB} (Management Information Base) è un dizionario o "mappa" 
    che descrive la struttura dei dati di un dispositivo. 
    L'\textbf{OID} (Object Identifier) è l'indirizzo numerico specifico all'interno di quella mappa
    (es. \texttt{1.3.6.1.4.1.39165.1.1.0}) che identifica un singolo valore, come il numero seriale.
    \item[API (Application Programming Interface):] un'interfaccia o "ponte" software che \\
    permette a due sistemi o applicazioni diverse (es. il nostro script Python e Zabbix) 
    di parlarsi, passarsi comandi e scambiarsi dati in modo sicuro e automatizzato.
    \item[ISAPI:] un set di API specifico e proprietario sviluppato da Hikvision.
    Permette di interrogare e comandare le telecamere tramite protocollo HTTP
    e di ricevere flussi continui di avvisi di sicurezza.
    % \item[Script / Middleware:] un programma (in questo progetto scritto in linguaggio Python)
    %  che funge da intermediario. Si occupa di tradurre, filtrare o trasportare dati da un punto A
    %   (es. file Excel o telecamera) a un punto B (Zabbix).
    \item[GIS (Geographic Information System):] sistemi informatici (come QGIS) 
    progettati per acquisire, memorizzare, analizzare e gestire dati georeferenziati.
    Nel caso di Zabbix permettono di visualizzare gli apparati guasti direttamente su una mappa cartografica.
    % \item[Payload:] in informatica, è il "carico utile" di un pacchetto di dati. 
    % Quando lo script aggiorna Zabbix, il payload è il blocco di testo (solitamente in formato JSON)
    %  che contiene le vere e proprie informazioni da modificare (es. nome e coordinate dell'host).
\end{description}
}

\chapter{L'Infrastruttura e il  monitoraggio della Smart City}

\section{Architettura generale dell'infrastruttura}
L'infrastruttura che necessita di monitoraggio nel progetto \emph{Smart City - Edison} 
è costituita da una serie di sottoreti virtuali \emph{VLAN} che si sviluppano sul territorio, 
comprendente circa 30 Comuni. In esse sono collegati diversi host, di natura eterogenea,
per oltre 1000 dispositivi attivi. 
Il sistema \textbf{Zabbix 7.0 LTS} è
ospitato su un ambiente server Linux (Ubuntu 24.04).

In precedenza, il monitoraggio delle sottoreti era affidato a Paessler PRTG Network Monitor.
La migrazione a Zabbix è stata una scelta dettata da fattori di costo dovuti alle licenze a numero limitato 
di host monitorabili di PRTG e da fattori di automatizzazione, le API di Zabbix sono 
più facilmente utilizzabili e meglio documentate, per la loro natura open-source.

Per ogni comune posso avere 4 sottoreti, per esempio nel comune di \lstinline|Lavenone|:
\begin{description}
    \item [TVCC-LAVE:] contiene host che sono telecamere e un gateway e un possibile server.
    \item [SW-LAVE: ]contiene host che sono switch e un gateway e un possibile server.
    \item [AP-LAVE: ]contiene host che sono access point e un gateway.
    \item [LTE-LAVE:] sottorete facoltativa che contiene host di ogni tipo che si basano su 
    tecnologia LTE o satellitare per comunicare.
\end{description}

La \textbf{funzionalità principale del sistema} consiste nell'applicare il template \emph{ICMP ping}
per ogni host presente nelle VLAN del progetto.
Questo template permette di applicare a tutti gli host un comando \emph{ping}, ogni minuto, e misurare,
tra 3 diversi item, l'accessibilità dell'host sulla rete. In caso di comunicazione interrotta 
il relativo trigger lancia un problem con \lstinline|Severity: high|.

\textbf{Warning: sarebbe bello inserire un diagramma di rete visivo.  
Un'immagine aiuterebbe a comprendere come Zabbix comunica con gli edge router e le telecamere.}

\subsection{VM Zabbix}
L'applicativo Zabbix viene eseguito su una virtual machine gestita da Nutanix, sui server Secoval, chiamata 
"SECZABBIX-TLC". Il sistema operativo utilizzato è Ubuntu Server (Ubuntu 24.04). Utente: "administrator".
Il sistema Zabbix è raggiungibile da http://192.168.100.57/zabbix con utente "Admin".
Zabbix si appoggia su un servizio MySql con utente: "zabbix@localhost".

\subsubsection{Servizi fondamentali per il funzionamento}
Questi sono i demoni e i servizi che devono rimanere sempre \inlinecode{active (running)}, per il funzionamento del sistema.
\begin{lstlisting}[language=bash, caption=Comandi utili per la verifica dell'operatività del sistema]
systemctl status zabbix-server #core
systemctl status apache2  #web server che ospita l'interfaccia grafica
systemctl status mysql #database dove viene salvato tutto
systemctl status hikvision-monitor #demone per la ricezione in tempo reale degli allarmi dalle telecamere
\end{lstlisting}

\subsubsection{Ambiente Python venv}
Ogni script python viene eseguito nell'ambiente virtualizzato e isolato \emph{venv}. 
In questo modo si mantiene un ambiente isolato e protetto. 
Qui sotto i comandi utili per interagire con venv e python:
\begin{lstlisting}[language=bash, caption=Comandi utili per venv e  l'ambiente python]
source venv/bin/activate #enter venv
deactivate #exit venv

#in venv:
pip list # lists all installed packages in the venv and their versions
pip install <package-name> # installs the package
pip install --upgrade <package-name> # upgrades the package to the latest available version
pip uninstall <package-name> # removes a package and its dependencies
python --version # shows which python version is currently active
\end{lstlisting}

\subsubsection{Stato dello storage}
Questo è lo stato del disco alla data 06/05/2026. Questa informazione può essere utile per valutare la crescita del sistema nel tempo.
\begin{lstlisting}[language=bash, caption=df\, disk free\, command: total\, used\, and available space for all mounted  filesystems]
administrator@seczabbix-tlc:~$ df -h
Filesystem                         Size  Used Avail Use% Mounted on
tmpfs                              1.6G  1.2M  1.6G   1% /run
/dev/mapper/ubuntu--vg-ubuntu--lv  193G   55G  129G  30% /
tmpfs                              7.8G     0  7.8G   0% /dev/shm
tmpfs                              5.0M     0  5.0M   0% /run/lock
/dev/sda2                          2.0G  200M  1.6G  11% /boot
tmpfs                              1.6G   16K  1.6G   1% /run/user/1000
\end{lstlisting}

\section{Le modalità del passaggio da PRTG a Zabbix}
Essendo le sottoreti da monitorare già presenti in PRTG, con il nome desiderato, è stato esportato 
un file \emph{json} contenente un array di oggetti così strutturati:
\begin{lstlisting}[language=c, caption=Struttura della sottorete estratta da PRTG]
    {
        "PRTG_ID":  4591,
        "Full_Path":  "smartCity/LAVENONE/TVCC-LAVE",
        "GroupName":  "TVCC-LAVE",
        "IPv4_Subnet":  "10.82.241.0/25"
    },
\end{lstlisting}
Questo file è stato utilizzato per inizializzare in Zabbix le stesse sottoreti, 
attraverso uno script python collegato tramite API: \lstinline|/home/administrator/usefulFilesForZabbix/import_zabbix.py|.

In Zabbix gli host possono essere catalogati in \emph{host groups}, 
a cui vengono assegnati da delle \emph{actions} chiamate \emph{discovery actions} 
dopo essere stati scoperti da delle \emph{discovery rules}: sistemi automatici 
che trovano gli host presenti in una determinata sottorete.

I compiti svolti da questo script sono stati:
\begin{itemize}
    \item Creazione di un nuovo \emph{host group} per ogni \lstinline|"Full_Path"|
    \item Creazione di una nuova \emph{discovery rule} per ogni \lstinline|"Full_Path"| 
    contenente la sottorete IPv4 e la modalità con cui può scoprire nuovi host, 
    in questo caso \lstinline|ICMP ping|
    \item Creazione di un nuovo \emph{discovery action} per ogni \lstinline|"Full_Path"| che assegni
    l'host appena scoperto dalla discovery rule al suo host group 
    e allo specifico template relativo alla tipologia di sottorete a cui appartiene
    (templateSC-TVCC, templateSC-SW, templateSC-AP, templateSC-LTE).
\end{itemize}

\section{Il funzionamento dei template ottenuti da "ICMP ping" sui gruppi di host}%da rivedere e completare
I template appena citati (templateSC-TVCC, templateSC-SW, templateSC-AP, \\templateSC-LTE) 
sono stati ottenuti clonando il template predefinito \emph{ICMP ping} e modificandolo.
In particolare l'unico item e relativo trigger mantenuti \emph{Enabled} nei template è \emph{ICMP ping}, 
la cui descrizione è: \\\emph{The host accessibility by ICMP ping.\\
0 - ICMP ping fails;\\
1 - ICMP ping successful.}

Il trigger relativo è descritto successivamente.

\chapter{Automazione dell'assegnazione dei metadati degli host - Inventory}

\section{Il problema}
Una volta scoperti gli host e memorizzati nel database di Zabbix essi sono quasi inutili, 
perché non contengono alcuna informazione che li riguarda. 
Molte delle loro caratteristiche, come nome e l'ID del palo a cui sono montati, 
sono presenti in fogli Excel di configurazione con molte altre informazioni.

Per collegare queste informazioni all'infrastruttura di monitoraggio,
abbiamo sviluppato una serie di script Python che sfruttano la libreria \texttt{pandas} 
per l'elaborazione massiva dei dati tabellari e la libreria \texttt{pyzabbix} 
per l'interazione asincrona e sicura con le API di Zabbix. 
L'autenticazione verso Zabbix avviene esclusivamente tramite API Token associato a un utente di
servizio con permessi ben definiti:

\begin{lstlisting}[language=Python, caption=Autenticazione tramite API Token]
from pyzabbix import ZabbixAPI

# Inizializzazione connessione API
zapi = ZabbixAPI("http://192.168.100.57/zabbix")
# Autenticazione con token 
zapi.login(api_token='API_TOKEN_SEGRETO_GENERATO_DA_ZABBIX')
\end{lstlisting}

\section{Implementazione dell'aggiornamento degli host}

Ogni host in Zabbix viene identificato con un attributo "name" che contiene l'IP con cui è stato scoperto. 
Inoltre possiede anche un attributo "visible name" che lo rappresenta nel sistema. L'obiettivo è andare a estrarre 
dal file Excel di configurazione l'informazione del loro nome, da inserire nell'attributo "visible name" 
e l'informazione sull'id del palo su cui sono montati, da inserire le campo "location" del loro attributo
"inventory" (ogni host possiede un ampio elenco di campi da compilare nell'attributo "inventory").

In un file separato viene associata all'IDpalo anche la coordinata geografica. 
Quindi per ogni host con uno specifico IDpalo gli si deve associare anche le coordinate geografiche.

Siccome la quasi totalità delle sottoreti al momento della scrittura di questa relazione 
presenta numerosi host con svariati problemi, che impediscono al sistema di scoprirli tramite
ping, risulta impossibile assegnare i loro metadati in una sola volta.
Quindi, mano a mano che essi andranno online, verranno scoperti dal sistema che dovrà associare 
a essi le informazioni contenute nei file Excel. Proprio come per gli host scoperti inizialmente.
Questo implica che l'assegnazione di questi dati deve essere fatta periodicamente, 
con un periodo congruo. Per questo ho scelto di utilizzare script python, eseguiti periodicamente 
dal servizio \texttt{crontab} presente nativamente in Ubuntu.

Nella pratica vengono utilizzati degli script python eseguiti nella macchina virtuale che 
ospita Zabbix. Tutti gli script ausiliari vengono salvati nella directory  
\texttt{/home/administrator/zabbixHostsScripts} (con alcune eccezioni) e tutte le esecuzioni sono svolte nell'ambiente
virtualizzato \texttt{venv} della stessa directory.

\subsection{Associazione host name e IDpalo}
%aggiungere comune e progetto edison
Per associare l'hostname e l'IDpalo allo specifico host, identificato dall'IP, viene utilizzato lo script:
\lstinline|zabbix_HostName_IDpalo_updater.py|.
Questo script utilizza \texttt{pandas} per la gestione del file .xlsx chiamato \lstinline|HostInfosUbuntu.xlsx|, che 
contiene il nome dell'host da inserire in "visible Name" e il nome del palo, "IDpalo", da inserire in 
"inventory>location".
Viene creato un dizionario degli host in Zabbix, comparati tramite IP con le informazioni presenti 
in Excel e l'host viene aggiornato solo se c'è incongruenza. 
Gli IP vengono normalizzati (si sistema il formato, ecc, ecc...) in una funzione separata(non è al 100\% perfetta,
in alcuni casi limite ci può essere ambiguità).

Questo script viene eseguito periodicamente tramite questo comando crontab, accessibile con il comando \inlinecode{crontab -e}: %inlineCode formatting
\begin{lstlisting}[language=Python, caption=setup of periodic execution in Crontab]
0 * * * * /home/administrator/zabbixHostsScripts/venv/bin/python3 /home/administrator/zabbixHostsScripts/zabbix_HostName_IDpalo_updater.py >> /home/administrator/zabbixHostsScripts/cron_log.log 2>&1
\end{lstlisting}

\subsection{Associazione coordinate, indirizzo e tecnologia}

Per associare le coordinate, l'indirizzo e la tecnologia a ogni host identificato dall'IDpalo
viene utilizzato lo script:
\inlinecode{zabbix_palo_location_and_info_updater.py}. 
Questo script utilizza \texttt{pandas} per la gestione del file .xlsx chiamato \\
\texttt{\detokenize{info_e_location_pali_ubuntu.xlsx}}, che 
contiene il nome del palo, "IDpalo", associato alle colonne 
"Tecnologia", "Indirizzo", "Dec. Latitudine" e "Dec. Longitudine" da inserire per ogni host in 
inventory>location nei vari campi: HW architecture, Site address A, Location latitude,
Location longitude. Le informazioni di "Dec. Latitudine" e "Dec. Longitudine" devono essere aggiornate 
ogni volta che si inseriscono nuovi possibili pali nel sistema monitorato.
Nello script viene creato un dizionario degli host in Zabbix, comparati tramite IDpalo con le informazioni presenti 
in Excel e l'host viene aggiornato solo se c'è incongruenza. 

Questo script viene eseguito periodicamente tramite questo comando crontab, accessibile con \inlinecode{crontab -e}:
\begin{lstlisting}[language=Python, caption=ciao]
19 * * * * /home/administrator/zabbixHostsScripts/venv/bin/python3 /home/administrator/zabbixHostsScripts/zabbix_palo_location_and_info_updater.py >> /home/administrator/zabbixHostsScripts/cron_log.log 2>&1
\end{lstlisting}

\subsection{Memorizzazione dei messaggi di log}
Gli script python visti fino a ora generano degli output testuali, che andrebbero nello standard output,
ma vengono reindirizzati dai job crontab in uno specifico file di log: 
\lstinline|cron_log.log|, ogni scrittura nel file è preceduta dal timestamp e dallo script che la genera.

Per evitare di memorizzare questi dati all'infinito e generare un file di testo gigantesco, 
ho utilizzato il sistema \textbf{Logrotate}, una utility Linux il cui scopo è la rotazione, 
compressione e rimozione di file di log automatici.

In questo caso il file contenente le configurazioni per il job logrotate del file \lstinline|cron_log.log| è accessibile tramite: \inlinecode{nano /etc/logrotate.d/zabbix_scripts}.
Al suo interno troviamo: 
\begin{lstlisting}[language=python, caption=Configurazione logrotate per \lstinline|cron_log.log|]
#path of logfile to manage
/home/administrator/zabbixHostsScripts/cron_log.log {
    #user privileges
    su administrator administrator
    #time period of rotation
    weekly
    #how many rotated and compressed files to keep
    rotate 12
    #add date to filename
    dateext
    dateformat -%Y-%m-%d
    #where to archive rotated logs
    olddir /home/administrator/zabbixHostsScripts/storico_log
    #compress the files but wait one cycle
    compress
    delaycompress
    #ignore errors if logfile is missing
    missingok
    #don't rotate if logfile is 0 bytes
    notifempty
    #create - permission - user ownership 
    create 0644 administrator administrator
}
\end{lstlisting}

\chapter{Estrazione Dati Host tramite Template Zabbix}

\section{Acquisizione informazioni sugli host}
Oltre alle informazioni riguardanti la posizione fisica dell'host 
è necessario associare ai singoli host informazioni sulla loro natura hardware e software.
Le modalità per acquisire queste informazioni sono varie e differenziate tra tipologia di host.
Questo è un limite; per ora solo la maggior parte degli host appartiene a una categoria abbastanza 
numerosa da giustificare un sistema automatico per acquisire queste informazioni.

I metodi utilizzati sono:
\begin{itemize}
    \item  \textbf{template con chiamata HTTP} per le telecamere di marca Hikvision.
    \item  \textbf{template con richiesta SNMP} per gli AP e gli SW.
\end{itemize}
Per ogni categoria: 
\begin{itemize}
    \item Creazione di un nuova \emph{discovery rule} per ogni sottorete coinvolta 
    (che possa contenere host da cui estrarre dati). Con determinati \emph{Checks} per filtrare preventivamente
    gli host.
    \item Creazione di un nuovo \emph{template} che possa comunicare e leggere dagli host le informazioni che espongono.
    \item Creazione di un nuova \emph{discovery action} per ogni categoria di acquisizione dati,
    che assegni l'host appena scoperto dalla discovery rule allo specifico template 
    relativo alla categoria di host per l'acquisizione dati a cui appartiene.
\end{itemize}

\section{Configurazione dell'host Inventory}
Siccome le informazioni estratte dai template devono essere memorizzate sugli host, 
in modo automatico, l'inventory degli host deve essere impostato sull'\texttt{Inventory mode} \textbf{Automatic}. 
Questo viene specificato nelle \emph{discovery action} nel menu \emph{operations} creando una operazione: \texttt{Set host inventory mode: Automatic}.

All'interno dei Template preconfigurati che sono stati adattati allo scopo, alcuni item
sono configurati per riversare il dato raccolto in uno specifico
campo dell'inventory tramite l'opzione \emph{Populates host inventory field},
mappando il valore su campi come \texttt{Model}, \texttt{Software} o \texttt{Serial number A}. 
I meccanismi esatti di estrazione variano in base alla tipologia e al brand dell'apparato.

\section{Template telecamere Hikvision}
Per gli host che sono telecamere hikvision (tutte modello: DS-2CD3666G2-IZS)
sono state create automaticamente le discovery rules (con script \texttt{/home/administrator\\/usefulFilesForZabbix/DiscoveryRuleCreationForHikvison\\/createDiscoveryRulesForHikvisions.py})\\
seguendo questa schema per il nome delle discovery rules, ad esempio: \\\lstinline|Discovery - smartCity_HikvisionCameras_AGNOSINE_LTE-AGNO|.

Esse implementano sulle sottoreti \textbf{TVCC} e \textbf{LTE}, identificate da subnet/subnetmask, 2 checks: 
collegamento HTTP port 80 e collegamento TCP port range 8443 (che solo queste telecamere espongono). 
In questo modo vengono filtrati gli host di altro tipo, 
come switch e server che risiedono sulle stesse sottoreti.

Queste discovery rules vengono eseguite ogni ora e le loro risposte vengono elaborate dalla discovery action 
\lstinline|smartCity_Auto-Configurazione_Hikvision| che applica il template e setta l'inventory mode su automatic.

Per ottenere il template ne è stato modificato uno dedicato (\lstinline|Hikvision camera by HTTP|)
ottenendo \lstinline|templateSC_configurazione_hikvision|, che interroga tramite API HTTP con autenticazione le telecamere,
e si occupa di leggere le informazioni e inserirle negli appropriati campi dell'inventory.
Lo username e password sono standard per tutte le telecamere, e sono contenuti in delle macro all'interno del template: {\$PASSWORD} e {\$USER}.

\section{Apparati standard (switch, access point, LTE)}
Per gli apparati che supportano facilmente l'utilizzo di SNMP come Switch e Access Point si è scelta
quella tecnologia, anche se non è possibile estrarre un set molto ampio di informazioni.

Per questi host sono state create automaticamente le discovery rules (con script \texttt{/home/administrator\\/usefulFilesForZabbix/DiscoveryRuleCreationForHikvison\\/apCreateDiscoveryRules.py e /swCreateDiscoveryRules.py})\\
seguendo questa schema per il nome delle relative discovery rules, ad esempio: \\\lstinline|Discovery - smartCity_APcambium_CAPOVALLE_AP-CAPO| 
e \lstinline|Discovery - smartCity_SWsenzaMarca_MAZZANO_SW-MAZZ|.

Esse implementano sulle sottoreti \textbf{SW} e \textbf{AP} un Check per filtrare gli host di altro tipo: 
collegamento SNMPv2 agent port 161 con SNMP community public e SNMP OID 1.3.6.1.2.1.1.1.0 (OID System Description, all'interno del MIB-II standard).

Queste discovery rules vengono eseguite ogni ora e le loro risposte vengono elaborate dalle discovery action 
\lstinline|smartCity_Auto-smartCity_Auto-Configurazione_APcambium| e \lstinline|smartCity_Auto-Configurazione_SWsenzaMarca| 
che applicano il template adatto e settano l'inventory mode su automatic.

Per ottenere i template ne è stato modificato uno dedicato (\lstinline|Generic by SNMP|) ottenendo
\lstinline|templateSC_Generic by SNMP_AP| per gli AP e \lstinline|templateSC_Generic by SNMP_SW| per gli SW. 
Essi interrogano tramite SNMP gli host, leggono le informazioni e le inseriscono negli appropriati campi dell'inventory.

Per l'esplorazione dei campi raggiungibili in SNMP ho usato questo strumento CLI:

\begin{lstlisting}[language=bash, caption=Esempio di discovery manuale dei MIB via CLI]
snmpwalk -v2c -c public 10.0.50.5 
\end{lstlisting}
%corretto fino a qui

\chapter{Sicurezza e rilevamento manomissioni (ISAPI)}

Essendo questo un progetto di gestione di telecamere di sicurezza abbastanza moderne è stato valutato come utile l'utilizzo dei sistemi di 
intelligenza artificiale implementati a bordo per rilevare diversi tipi di eventi che accadono alla singola telecamera.
Tra questi troviamo la possibilità di rilevare l'oscuramento dell'obiettivo della telecamera. 
Per sfruttare questa funzionalità è stato creato un sistema di raccoglimento e invio automatico di avvisi, integrandolo con Zabbix.

\section{Descrizione dell'implementazione}
Il sistema gestisce 2 eventi provenienti dalle telecamere: il  tamperdetection o l'oscuramento dell'obiettivo e il motiondetection o il rilevamento di movimento. 
Il secondo è stato implementato solo per motivi di debug, dato che condivide gli stessi canali dell'oscuramento ma è molto più facilmente riproducibile.

Il workflow che viene eseguito per garantire la corretta gestione di queste informazioni si divide in questi step:
\begin{itemize}
    \item Estrazione degli IP delle telecamere Hikvision presenti in Zabbix con uno script python. 
    \item Ogni IP estratto è utilizzato per lanciare verso la telecamera le configurazioni dei loro servizi tamperdetection e motiondetection tramite configurazioni in formato XML.
    \item Si apre una connessione http con ogni telecamera mantenuta perpetuamente aperta in modo asincrono, attraverso uno script python. Su queste connessioni le telecamere comunicano i loro avvisi
    ed essi vengono comunicati a Zabbix sugli appositi Item dei trigger tramite un API CLI di Zabbix.
    \item All'interno del template relativo alle telecamere Hikvision abbiamo 2 trigger per motiondetection e oscuramento con le loro impostazioni.
    \item I log di questi script vengono gestiti da logrotate salvandoli fino a 3 mesi.
\end{itemize}
\textbf{Bottleneck}: se molti allarmi si attivano contemporaneamente, 
lo script python lancia altrettanti script bash e ciò sovraccaricherebbe la CPU. 
Se l'utilizzo rimane per il tamperdetection e solo saltuariamente per motiondetection non è un problema. 

\section{Estrazione degli IP delle telecamere Hikvision}

È stato creato uno script python in \\\inlinecode{/home/administrator/zabbixHostsScripts/hikvisionOscuramentoScripts}\\\inlinecode{/estraiIPtelecamereHikvision.py}
che si collega tramite API token a Zabbix, estrae tutti gli host, con il relativo campo "contact" dell'inventory, per filtrarli, e salva l'indirizzo IP solo
se nel "contact" abbiamo una stringa corrispondente a "Hikvision.China": \\
\inlinecode{if inventario.get("contact") == "Hikvision.China":}
Gli IP sono salvati nel file \inlinecode{ip_hikvision.txt} nella stessa directory, sovrascrivendone il contenuto per aggiornare la lista. 

Questo script viene eseguito periodicamente ogni ora dal servizio crontab con un \textit{cronjob} 
sotto l'utente \texttt{administrator} e salva lo standard output in un file di log: 
\\\inlinecode{/home/administrator/zabbixHostsScripts/hikvisionOscuramentoScripts}\\\inlinecode{/ipHik_log.log}

\begin{lstlisting}[language=bash, caption=Automazione crontab estrazione degli IP delle telecamere Hikvision]
#Esecuzione oraria al minuto 35
35 * * * * /home/administrator/zabbixHostsScripts/venv/bin/python3 /home/administrator/zabbixHostsScripts/hikvisionOscuramentoScripts/estraiIPtelecamereHikvision.py >> /home/administrator/zabbixHostsScripts/hikvisionOscuramentoScripts/ipHik_log.log 2>&1
\end{lstlisting}

\section{Configurazioni tamperdetection e motiondetection hikvision tramite XML}

Le telecamere Hikvision espongono una API HTTP che permette di modificare le configurazioni e
le impostazioni della telecamera tramite l'upload di file semi strutturati XML. 
Essi riguardano il setup delle impostazioni e della modalità di comunicazione per 2 eventi: tamperdetection e motiondetection 
(mai utilizzati automaticamente, sono commentati). 

Questi file si trovano in 
\\\inlinecode{/home/administrator/zabbixHostsScripts/hikvisionOscuramentoScripts/XMLs}\\ e quelli riguardanti il tamperdetection/oscuramento sono:
\inlinecode{oscuramento_setup.xml} e \inlinecode{triggers_oscuramento_setup.xml}. I quali rispettivamente gestiscono l'attivazione e l'impostazione dell'area dell'immagine in cui la telecamera è 
sensibile all'oscuramento e la comunicazione dell'evento triggerato all'\texttt{EventTriggerNotification center} (lancia la comunicazione http).

\subsection{Implementazione del sistema per ogni telecamera}
Per ogni IP individuato dallo script visto in precedenza, chiamiamo 2 API, tramite uno script bash: \\
\inlinecode{/home/administrator/zabbixHostsScripts/hikvisionOscuramentoScripts/}\\
\inlinecode{configura_telecamere.sh}.
Questo script utilizza la utility, per CLI, \emph{curl} (Client URL) con metodo HTTP PUT
per caricare sulla telecamera con autenticazione le configurazioni XML viste prima.

\begin{lstlisting}[language=bash, caption=Chiamata API hikvision per configurare il tamperdetection configura\_telecamere.sh]
  # --- 1. Attivazione sensore Tampering ---
  #CURL FLAGS
  #-s	--silent        no graphical output
  #--digest             digest auth, hashed password
  #-u	--user          credentials
  #-X PUT	--request   HTTP method
  #-d	--data          data to be sent 
  #-H	--header        header, XML content
  RESPONSE=$(curl -s --digest -u "$USER:$PASS" -X PUT -d "@XMLs/oscuramento_setup.xml" "http://$ip/ISAPI/System/Video/inputs/channels/1/tamperDetection" -H "Content-Type: application/xml")
  #grep flag -q is for quite, no output, only 1 for found 0 for not found
  if echo "$RESPONSE" | grep -q "<statusCode>1</statusCode>"; then
    echo " [-] Setup Tampering: OK"
  else
    echo " [X] Setup Tampering: ERRORE"
    host_has_errors=1
  fi

  # --- 2. Trigger Tampering al Centro ---
  RESPONSE=$(curl -s --digest -u "$USER:$PASS" -X PUT -d "@XMLs/triggers_oscuramento_setup.xml" "http://$ip/ISAPI/Event/triggers/tamper-1" -H "Content-Type: application/xml")
  if echo "$RESPONSE" | grep -q "<statusCode>1</statusCode>"; then
    echo " [-] Trigger Tampering: OK"
  else
    echo " [X] Trigger Tampering: ERRORE"
    host_has_errors=1
  fi
\end{lstlisting}

Infine, dopo l'esecuzione viene stampato un riepilogo: 

\begin{lstlisting}[language=bash, caption=Riepilogo configura\_telecamere.sh]

echo "================================================================="
echo "RIEPILOGO ESECUZIONE"
echo "================================================================="
echo "Terminato il: $(date '+%Y-%m-%d %H:%M:%S')"
echo "Totale telecamere processate: $TOTAL_HOSTS"
echo ""
echo "=> Configurate con successo (tutto OK): $SUCCESS_HOSTS su $TOTAL_HOSTS"
echo "=> Fallite (almeno un errore): $ERROR_HOSTS su $TOTAL_HOSTS"
echo "================================================================="
\end{lstlisting}

\subsection{Cronjob e log}
Questo script viene eseguito periodicamente ogni ora dal servizio crontab con un \textit{cronjob} 
sotto l'utente \texttt{administrator} e salva lo standard output in un file di log: 
\\\inlinecode{/home/administrator/zabbixHostsScripts/hikvisionOscuramentoScripts}\\\inlinecode{/config_telecamere_log.log}

\begin{lstlisting}[language=bash, caption=Automazione crontab configurazione dell'oscuramento delle telecamere Hikvision]
#Esecuzione oraria al minuto 2
2 * * * * /home/administrator/zabbixHostsScripts/hikvisionOscuramentoScripts/configura_telecamere.sh >> /home/administrator/zabbixHostsScripts/hikvisionOscuramentoScripts/config_telecamere_log.log 2>&1
\end{lstlisting}

\section{Ricezione degli eventi e invio a Zabbix}
È stato creato uno script python in \\\inlinecode{/home/administrator/zabbixHostsScripts/hikvisionOscuramentoScripts}\\\inlinecode{/hikvision_listener_multiplehosts.py}
che si comporta come middleware. Il suo scopo principale è fare da ponte tra le telecamere di videosorveglianza Hikvision e il server di monitoraggio Zabbix, 
intercettando gli allarmi in tempo reale e convertendoli valori di item nel sistema Zabbix.
 
Invece di interrogare continuamente le telecamere (polling),
lo script si mette in "ascolto" continuo sui loro flussi di eventi (streaming), 
catturando istantaneamente anomalie come il rilevamento del movimento (VMD) 
o l'oscuramento della telecamera (Tampering/Shelter alarm).

\subsection{Implementazione del listener}
Lo script si collega tramite API token a Zabbix, utilizzando la libreria pyzabbix, 
estrae tutti gli host, con il relativo campo "contact" dell'inventory, per filtrarli, e salva l'indirizzo IP solo
se nel "contact" abbiamo una stringa corrispondente a "Hikvision.China".
Questo rende il sistema altamente scalabile: se aggiungi una nuova telecamera su Zabbix, 
il riavvio successivo dello script la includerà automaticamente.

Per ogni telecamera trovata, viene avviato un "task" indipendente grazie alla libreria asyncio(Funzione ascolta\_telecamera). 
Questo permette di gestire decine di telecamere contemporaneamente usando un solo processo leggero, 
senza che l'attesa di dati da una telecamera blocchi le altre.
Lo script utilizza http per aprire una connessione HTTP persistente (Stream) verso l'endpoint ISAPI delle telecamere (/ISAPI/Event/notification/alertStream).
In questo flusso continuo, la telecamera invia frammenti in formato XML ogni volta che accade qualcosa.
Lo script legge il flusso riga per riga, isola i blocchi XML (<EventNotificationAlert>) e li analizza in tempo reale.
Estrae il tipo di evento (eventType) e il suo stato (eventState, se attivo o rientrato).
Quando viene rilevato un evento valido (es. movimento o telecamera coperta), lo script deve avvisare Zabbix.
Per farlo lo script esegue un comando di sistema: zabbix\_sender. 
Questo è uno strumento a riga di comando nativo di Zabbix, 
progettato specificamente per inviare dati ad alte prestazioni a Zabbix.
Il comando comunica l'hostname esatto della telecamera, la chiave dell'allarme 
(es. hikvision.motionDetection) e il valore (1 per allarme attivo, 0 per allarme rientrato).

\subsubsection{Items in Zabbix}
Nel template \texttt{templateSC\_configurazione\_hikvision} sono stati creati 2 nuovi Item di tipo \emph{Zabbix trapper}: 
AllarmeItemOscuramentoProva(key: hikvision.oscuration) e AllarmeItemMotionDetectionProva(hikvision.motionDetection). %valutare se togliere prova
Essi sono aggiornati solo da zabbix\_sender. 

\subsection{Esecuzione continuativa e log}

Lo script viene configurato come servizio systemd, garantendo che si avvii automaticamente all'accensione del server.
Il file con le configurazioni è \\\inlinecode{/etc/systemd/system/hikvision-monitor.service}.
Se lo script dovesse bloccarsi per un errore critico, il sistema operativo lo riavvierà automaticamente entro 10 secondi.
È stata inserita una regola che "uccide" e riavvia il servizio ogni 24 ore in modo controllato. 
Questo è un ottimo accorgimento tecnico per evitare accumuli di memoria 
dovuti alle connessioni in streaming prolungate e per forzare il sistema a riscaricare la lista aggiornata delle telecamere da Zabbix ogni giorno.
Tutto l'output standard viene reindirizzato in un file di log \inlinecode{hikvision_listener_log.log} su disco per permettere il monitoraggio del corretto funzionamento.


\begin{lstlisting}[language=bash, caption=Configurazione Systemd per hikvision\_listener\_multiplehosts.py]
[Unit]
Description=Monitoraggio Eventi Hikvision (Tampering e VMD) verso Zabbix
After=network.target

[Service]
# Esegue Python in modalità "unbuffered" (-u) per scrivere subito i print nel log
ExecStart=/home/administrator/zabbixHostsScripts/venv/bin/python3 -u /home/administrator/zabbixHostsScripts/hikvisionOscuramentoScripts/hikvision_listener_multiplehosts.py
WorkingDirectory=/home/administrator/zabbixHostsScripts/hikvisionOscuramentoScripts

# Politiche di riavvio in caso di crash
Restart=always
RestartSec=10

# metodo per l'aggiornamento quotidiano:
# Forza il riavvio del servizio ogni 86400 secondi (24 ore)
RuntimeMaxSec=86400

# Gestione dei Log
StandardOutput=append:/home/administrator/zabbixHostsScripts/hikvisionOscuramentoScripts/hikvision_listener_log.log
StandardError=inherit

[Install]
WantedBy=multi-user.target

#Per controllare se lo script sta girando correttamente, chiedere lo stato a systemd:
#sudo systemctl status hikvision-monitor.service

#E per vedere l'output in tempo reale, usare il comando tail:
#tail -f /home/administrator/zabbixHostsScripts/hikvisionOscuramentoScripts/hikvision_listener_log.log
\end{lstlisting}

\section{Gestione dei log relativi al rilevamento manomissioni}
Gli script visti fino a ora generano degli output testuali, che andrebbero
nello standard output, ma vengono reindirizzati dai job crontab e systemd in specifici file di log: 
ipHik\_log.log, config\_telecamere\_log.log e hikvision\_listener\_log.log. 
Ogni scrittura nei file è preceduta dal timestamp.
Come visto in precedenza per evitare di memorizzare questi dati all'infinito e generare un file di testo
gigantesco, ho utilizzato il sistema Logrotate, una utility linux il cui scopo è la
rotazione, compressione e rimozione di file di log automatici.
In questo caso il file contenente le configurazioni per il job logrotate 
é accessibile tramite: \inlinecode{nano /etc/logrotate.d/hikvision_scripts}. Al
suo interno troviamo:

\begin{lstlisting}[language=bash, caption=Cronjob log relativi al rilevamento manomissioni]
# Blocco 1: Per gli script periodici (crontab)
/home/administrator/zabbixHostsScripts/hikvisionOscuramentoScripts/config_telecamere_log.log
/home/administrator/zabbixHostsScripts/hikvisionOscuramentoScripts/ipHik_log.log {
    su administrator administrator
    weekly
    maxsize 30M
    rotate 12
    dateext
    dateformat -%Y-%m-%d
    olddir /home/administrator/zabbixHostsScripts/hikvisionOscuramentoScripts/storico_log
    compress
    delaycompress
    missingok
    notifempty
    create 0644 administrator administrator
}

# Blocco 2: Per lo script Python in ascolto continuo (systemd)
/home/administrator/zabbixHostsScripts/hikvisionOscuramentoScripts/hikvision_listener_log.log {
    su administrator administrator
    weekly
    maxsize 15M
    rotate 12
    dateext
    dateformat -%Y-%m-%d
    olddir /home/administrator/zabbixHostsScripts/hikvisionOscuramentoScripts/storico_log
    compress
    delaycompress
    missingok
    notifempty
    copytruncate
}
\end{lstlisting}

\chapter{Trigger e Report Giornalieri}

Fino a ora si è parlato di come il sistema comunica le informazioni agli \emph{item} di Zabbix, ossia i sensori.
Adesso verranno descritti i \emph{trigger}, ossia le soglie, le regole logiche, che vengono eseguite sui dati forniti dagli item. 
Zabbix e i suoi template ne implementano molti, spesso con poco rilevanti per le informazioni che vogliamo estrarre dalla rete, 
quindi verranno descritti solo i \emph{trigger} che sono stati lasciati attivi, \emph{Enabled}, non quelli disattivati. 

Questi trigger sollevano \emph{problems}, relativi alla problematica che rilevano, e spesso, lanciano delle allerte tramite e-mail.
Oltre all'invio di queste notifiche per le criticità,
è emersa la necessità di fornire alla direzione e al team tecnico un riepilogo giornaliero. 
Poiché Zabbix non supporta nativamente l'invio programmato di file arbitrari o
analisi testuali complesse via email, è stata progettata un'architettura in tre fasi: 
esportazione locale in formato CSV dei \emph{problems} tramite \textit{Custom AlertScript}, elaborazione analitica tramite Python e invio automatizzato via \textit{Crontab}.

\section{Tutti i template e i relativi trigger}
\begin{enumerate} 
    \item \textbf{templateSC-AP} clonato da \emph{ICMP ping}
        \begin{itemize}
            %\item \textbf{ICMP ping} The host accessibility by ICMP ping
            \item \textbf{ICMP Ping: Unavailable by ICMP ping} \\
            \inlinecode{max(/templateSC-SW/icmpping, 1h)=0} \\
            Nell'ultima ora i valori dell'item "icmpping" sono tutti 0, ping fallito.
        \end{itemize}
    \item \textbf{templateSC-LTE} clonato da \emph{ICMP ping}
        \begin{itemize}
            \item \textbf{ICMP Ping: Unavailable by ICMP ping} \\
            \inlinecode{max(/templateSC-SW/icmpping, 1h)=0} \\
            Nell'ultima ora i valori dell'item "icmpping" sono tutti 0, ping fallito.
        \end{itemize}
    \item \textbf{templateSC-SW} clonato da \emph{ICMP ping}
        \begin{itemize}
            \item \textbf{ICMP Ping: Unavailable by ICMP ping} \\
            \inlinecode{max(/templateSC-SW/icmpping, 1h)=0} \\
            Nell'ultima ora i valori dell'item "icmpping" sono tutti 0, ping fallito.
        \end{itemize}
    \item \textbf{templateSC-TVCC} clonato da \emph{ICMP ping}
        \begin{itemize}
            \item \textbf{ICMP Ping: Unavailable by ICMP ping} \\
            \inlinecode{max(/templateSC-SW/icmpping, 1h)=0} \\
            Nell'ultima ora i valori dell'item "icmpping" sono tutti 0, ping fallito.
        \end{itemize}
    \item \textbf{templateSC\_configurazione\_hikvision} clonato da \emph{Hikvision camera by HTTP}
        \begin{itemize}
            \item \textbf{AllarmeMotionDetectionTrigger} \\
            \inlinecode{last(/templateSC_configurazione_hikvision/hikvision.motionDetection)=1 and nodata(/templateSC_configurazione_hikvision/hikvision.motionDetection,20s)=0} \\
            L'ultimo valore ricevuto dall'item "hikvision.motionDetection" è 1, ossia movimento percepito, e l'item ha ricevuto valori negli ultimi 20 secondi. 
            Questo trigger è Disabled, ma è presente per motivi di debug.
            \item \textbf{AllarmeOscurationTrigger} \\
            \inlinecode{count(/templateSC_configurazione_hikvision/hikvision.oscuration, 30m, "eq", "1") >= 3} \\
            Si attiva solo se negli ultimi 30 min sono stati registrati almeno 3 valori di "hikvision.oscuration" a 1.
        \end{itemize}
    \item \textbf{templateSC\_Generic by SNMP\_AP} clonato da \emph{Generic by SNMP}
        \begin{itemize}
            \item Nessun trigger Enabled, solo gli item per l'estrazione di dati tramite SNMP.
        \end{itemize}
    \item \textbf{templateSC\_Generic by SNMP\_SW} clonato da \emph{Generic by SNMP}
        \begin{itemize}
            \item Nessun trigger Enabled, solo gli item per l'estrazione di dati tramite SNMP.
        \end{itemize}
    \item \textbf{templateFuoriProgettoEdison\_Barghe\_ICMP Ping} clonato da \emph{ICMP ping}
        \begin{itemize}
            \item \textbf{ICMP Ping: Unavailable by ICMP ping} \\
            \inlinecode{max(/templateSC-SW/icmpping, 1h)=0} \\
            Nell'ultima ora i valori dell'item "icmpping" sono tutti 0, ping fallito.
        \end{itemize}
\end{enumerate}

\section{Georeferenziazione  e integrazione GIS integrato in Zabbix}
I metadati raccolti in Zabbix, come longitudine e latitudine, vengono utilizzati nativamente dal software GIS integrato (geographic information system).
Le coordinate geografiche degli apparati vengono estratte dai campi \texttt{location\_lat} (Latitudine) e \texttt{location\_lon} (Longitudine) dell'inventory degli host.
All'interno della dashboard \emph{Mappa}, si possono visualizzare gli host su di una cartina, codificati tramite colore per la gravità dei problemi che eventualmente riscontrano.

\section{Alerting tramite email}
I trigger visti fino ad ora lanciano dei \emph{Problem}, spesso di \emph{Severity: High}.
Essi sono memorizzati e gestiti da Zabbix internamente, ma è utile, per Secoval, avere degli alert via mail che contengano i dettagli dei problemi.

\subsection{Setup delle email}
\begin{itemize}
    \item Rendere \emph{Enabled} il media type \textbf{Email} e inserire le informazioni del server di posta elettronica a cui si aggancia. \\
        Alerts>Media types>Email;
    \item Set up dello user group \emph{EmailAlerts}
    \item Set up degli user, nome e pw, media(type: email, indirizzo dell'utente a cui si vuole che la mail arrivi)
\end{itemize}

\subsection{Setup delle trigger actions}
\begin{tabular}{|p{3cm}|p{5.25cm}|p{5.25cm}|}
\hline
\textbf{Name} & \textbf{Conditions} & \textbf{Operations} \\
\hline
% Row 1
oscuramento & 
Trigger equals templateSC\_configurazio\-ne\_hikvision: AllarmeOscurationTrigger & 
Send message to users: Brando Calderara, CED\_secoval via Email 
\newline
Send message to users: Admin (Zabbix Administrator) via CustomExportProblemsToCSV \\
\hline
% Row 2
motionDetection & 
Trigger equals templateSC\_configurazio\-ne\_hikvision: AllarmeMotionDetectionTrigger & 
Send message to users: Brando Calderara via Email 
\newline 
Send message to users: Admin (Zabbix Administrator) via CustomExportProblemsToCSV \\
\hline
% Row 3
Allarme Ping Down Unavailable - Progetto Edison SmartCity & 
Trigger equals templateSC-TVCC: ICMP Ping: Unavailable by ICMP ping\newline
Trigger equals templateSC-SW: ICMP Ping: Unavailable by ICMP ping\newline
Trigger equals templateSC-AP: ICMP Ping: Unavailable by ICMP ping\newline
Trigger equals templateSC-LTE: ICMP Ping: Unavailable by ICMP ping & 
Send message to users: Brando Calderara via Email 
\newline 
Send message to users: Brando Calderara, CED\_secoval via Email\newline
Send message to users: Admin (Zabbix Administrator) via CustomExportProblemsToCSV \\
\hline
% Row 4
Allarme Ping Down Unavailable - FuoriProgettoEdison/Barghe & 
Trigger equals templateFuoriProgettoEdison\_Barghe\_ICMP Ping: ICMP Ping: Unavailable by ICMP ping\newline
Host group equals FuoriProgettoEdison/Barghe & 
Send message to users: Brando Calderara via Email 
\newline 
Send message to users: Brando Calderara, Catalin Surdu, CED\_secoval, Sergio Scotuzzi via Email\newline
Send message to users: Admin (Zabbix Administrator) via CustomExportProblemsToCSV \\
\hline
\end{tabular}

Al loro interno, nelle trigger actions, in \emph{operations}, si aggiunge una operazione, di tipo Email, agli user interessati.
Sotto \emph{Custom Message} possiamo personalizzare l'email inviata, aggiungendo con le macro le informazioni utili sull'host e il problema. 
Stessa cosa in \emph{Recovery operations}.

\emph{Send message to users: Admin (Zabbix Administrator) via CustomExportProblemsToCSV} questa operazione verrà approfondita successivamente.

\section{Esportazione dati tramite custom AlertScript}
Per tracciare l'intero ciclo di vita di un problema (dalla sua apertura alla sua risoluzione), 
è stato creato un nuovo \textit{Media Type} di tipo \texttt{Script} in Zabbix. 
Ogni volta che un trigger si attiva o si ripristina, l'Action associata "invia" le 
informazioni a uno script Bash locale sul server Linux, anziché spedire una email.

Lo script Bash si occupa di formattare le Macro native di Zabbix e accodarle in un file 
\texttt{/var/log/zabbix/ExportedProblemsToCSV.csv}.

\begin{lstlisting}[language=bash, caption=Custom AlertScript per l'esportazione in CSV (/usr/lib/zabbix/alertscripts/ScriptCustomExportProblemsToCSV.sh)]
#!/bin/bash

# Define where to save the CSV file
LOGFILE="/var/log/zabbix/ExportedProblemsToCSV.csv"
#LOGFILE="/home/administrator/zabbixHostsScripts/dailyRecordsScripts/ExportedProblemsToCSV.csv"

# If the file doesn't exist yet, create it and write the Header row
if [ ! -f "$LOGFILE" ]; then
    echo "Status,EventID,Date,Time,Host,IP,Severity,EventName,TriggerName,RecoveryDate,RecoveryTime,Duration" > "$LOGFILE"
    # Make sure Zabbix has permission to write to it
    #chown zabbix:zabbix "$LOGFILE"
fi

# Append the incoming Zabbix data as a new row in the CSV
# Notice the curly braces for variables 10, 11, and 12!
echo "\"$1\",\"$2\",\"$3\",\"$4\",\"$5\",\"$6\",\"$7\",\"$8\",\"$9\",\"${10}\",\"${11}\",\"${12}\"" >> "$LOGFILE""
\end{lstlisting}

I parametri passati allo script dall'interfaccia web di Zabbix corrispondono a specifiche Macro, 
tra cui \texttt{{EVENT.STATUS}} (per distinguere \texttt{PROBLEM} da \texttt{RESOLVED}), \texttt{{EVENT.ID}}, e \texttt{{EVENT.RECOVERY.TIME}}. 

\section{Elaborazione analitica tramite Python}
Per dare significato ai dati grezzi estratti nel CSV, sono stati sviluppati degli script Python 
(\texttt{emailAndClearCsvProblemFile.py}, monthlyReport.py) che fungono da motore di analisi. 
Sfruttando la libreria nativa \texttt{csv}, nel primo caso, 
lo script legge tutte le righe registrate nelle 24 ore precedenti e le categorizza applicando specifiche logiche. 
Nel secondo caso si genera un report, che contiene diversi KPI.

\subsection{Categorizzazione degli eventi: emailAndClearCsvProblemFile.py}
Lo script Python elabora il file e divide le informazioni in quattro sezioni principali, che andranno a comporre il corpo dell'email:
\begin{itemize}
\item \textbf{Riepilogo:} Conteggio totale dei problemi rilevati e di quelli risolti automaticamente
 (riconciliando le righe tramite la combinazione di \texttt{Host} ed \texttt{EventID}).
\item \textbf{Problemi Non Risolti:} Una lista degli host che hanno registrato uno stato 
\texttt{PROBLEM} senza un corrispondente \texttt{RESOLVED} nello stesso intervallo di tempo.
\item \textbf{Infrastruttura Instabile (Oscillazione/Flapping):} Il programma conta quante volte un singolo apparato 
ha attivato l'allarme \textit{ICMP Ping: Unavailable}. 
Se il conteggio supera una soglia definita (es. 2 disconnessioni in un giorno), 
l'host viene flaggato come "instabile", suggerendo problemi di alimentazione o di segnale LTE anziché un semplice riavvio.
\item \textbf{Allarmi di Sicurezza Critici (Oscuramento):} Lo script filtra il CSV alla ricerca di eventi con 
\texttt{TriggerName} pari a \textit{AllarmeOscurationTrigger}. 
Se presenti, vengono messi in risalto per richiedere un'ispezione fisica immediata.
\end{itemize}

\begin{lstlisting}[language=Python, caption=Estratto del motore di analisi Python per la rilevazione delle oscillazioni]
Logica di scansione CSV e aggregazione

for row in reader:
status = row.get('Status', '')
host = row.get('Host', '')
trigger = row.get('TriggerName', '')

unique_key = f"{host}_{trigger}"

if status == "PROBLEM":
    total_problems += 1
    unsolved_dict[unique_key] = row
    
    # Conteggio cadute ICMP per rilevamento Flapping
    if "ICMP Ping" in trigger:
        icmp_counts[host] += 1
        
elif status == "RESOLVED":
    total_resolved += 1
    # Rimuove il problema dalla lista dei non risolti
    if unique_key in unsolved_dict:
        del unsolved_dict[unique_key]

Identifica host instabili (3 o piu cadute)

flapping_hosts = {host: count for host, count in icmp_counts.items() if count >= 3}
\end{lstlisting}

\section{Funzionamento: monthlyReport.py}
% \section{Automazione della Reportistica Mensile e Calcolo dei KPI}
% \label{sec:report_mensile}

Per garantire una visione d'insieme sullo stato di salute dell'infrastruttura, 
è stato sviluppato uno script Python dedicato alla generazione e all'invio di un report mensile. 
Questo strumento elabora eventi registrati da Zabbix, estrapolando metriche chiave (KPI).

\subsection{Logica di funzionamento}
Lo script agisce elaborando un archivio di file CSV generati quotidianamente, 
i quali contengono lo storico degli eventi (trigger) catturati da Zabbix. 
Attraverso la libreria \texttt{datetime}, il sistema itera sulle righe dei log filtrando e isolando 
esclusivamente gli eventi verificatisi nel mese solare precedente a quello di esecuzione. 
Questo approccio algoritmico garantisce che il report prodotto contenga dati consolidati e definitivi del mese appena trascorso. 

Inoltre, per contestualizzare i dati rispetto alle dimensioni attuali della rete, 
lo script interroga in tempo reale le API JSON-RPC di Zabbix tramite il metodo \texttt{host.get}. 
Questa chiamata, ottimizzata tramite il parametro \texttt{countOutput}, restituisce istantaneamente il numero esatto di host attivi 
e monitorati nel momento della generazione del report.

\subsection{Metriche e KPI calcolati}
L'analisi aggregata dei log permette di estrapolare diversi (KPI) per valutare la stabilità dell'infrastruttura Smart City:
\begin{itemize}
    \item \textbf{Volume degli eventi:} conteggio totale dei trigger entrati in stato \texttt{PROBLEM} 
    e di quelli risolti (\texttt{RESOLVED}) durante l'intero mese.
    \item \textbf{Analisi delle sottoreti:} suddivisione percentuale e numerica degli eventi problematici per tipologia di apparato,
     calcolata analizzando il terzo e quarto ottetto degli indirizzi IP (es. TVCC, Switch, Access Point, connettività LTE).
    \item \textbf{Classifica host critici:} identificazione dei 10 dispositivi (\textit{Top 10}) 
    che hanno generato il maggior numero di anomalie, permettendo di isolare i colli di bottiglia hardware o di rete.
    \item \textbf{Distribuzione temporale:} mappatura degli eventi per fascia oraria e distribuzione settimanale. 
    Questa metrica è cruciale per identificare pattern di disservizio legati a momenti specifici della giornata.
\end{itemize}

\subsection{Esecuzione pianificata e salvataggio dei log}
L'esecuzione dello script è completamente automatizzata e slegata dall'intervento umano grazie al demone \texttt{cron} di Ubuntu Server. 
È stato configurato un \textit{cron job} specifico per l'avvio alle ore 01:00 del primo giorno di ogni mese solare. 

Per garantire la tracciabilità delle esecuzioni e facilitare le operazioni di \textit{troubleshooting}, 
l'output standard e gli eventuali log di errore generati da Python vengono reindirizzati e accodati in un file di log locale 
(\texttt{cron\_report.log}) utilizzando la sintassi di ridirezione della shell (\texttt{>> \dots 2>\&1}).

\subsection{Notifica automatica via email}
Al termine della computazione, il report testuale formattato viene inviato in automatico ai referenti tecnici del CED aziendale.
Questo modulo sfrutta la libreria \texttt{smtplib} e \texttt{email.message} di Python per interfacciarsi con il server SMTP interno (\texttt{pmg.secoval.it}). 

% Sfruttando la configurazione di \textit{anonymous relay} intrinseca alle reti locali protette, 
% lo script stabilisce una connessione TCP sulla porta 25 (originata dall'indirizzo IP fidato del server Zabbix) e 
% provvede all'inoltro del messaggio senza la necessità di trasmettere credenziali di autenticazione in chiaro, 
% minimizzando la superficie di attacco e la complessità di gestione dei segreti.

\section{Automazione e gestione del ciclo di vita: emailAndClearCsvProblemFile.py}
L'ultimo tassello dell'architettura riguarda l'automazione dell'invio e la manutenzione dello spazio disco. 
Lo script Python utilizza la libreria \texttt{smtplib} per compilare il testo analizzato, 
allegare il file CSV per completezza di archiviazione, e instradare l'email attraverso il server SMTP interno 
(\texttt{pmg.secoval.it} sulla porta 25, senza necessità di autenticazione TLS grazie alle regole di relay interne).

Per garantire che l'esecuzione avvenga puntualmente ogni mattina,
e soprattutto per permettere allo script di avere i privilegi necessari per eliminare il file CSV 
(creato dall'utente di sistema \texttt{zabbix}), è stato configurato un \textit{cronjob} sotto l'utente \texttt{root}.

\begin{lstlisting}[language=bash, caption=Automazione crontab (eseguito come root)]
#Esecuzione giornaliera alle 06:10 del mattino

10 6 * * * /home/administrator/zabbixHostsScripts/venv/bin/python3 /home/administrator/zabbixHostsScripts/dailyRecordsScripts/emailAndClearCsvProblemFile.py >> /var/log/zabbix\_csv\_mailer.log 2>&1
\end{lstlisting}

Al termine dell'invio con esito positivo, lo script invoca \\\inlinecode{shutil.move(CSV_FILE, percorso_destinazione)}
posizionando il file nella directory di archivio. Questo approccio previene l'accumulo di vecchi alert,
garantendo che al verificarsi del primo errore della giornata successiva, 
l'AlertScript Bash rigeneri un file CSV "pulito", comprensivo di riga di intestazione.

\begin{lstlisting}[language=python, caption=Gestione dei dati in exportedProblemsToCSV.csv raccolti da Zabbix]
# ==========================================
# 5. GESTIONE FILE (ARCHIVIAZIONE)
# ==========================================
try:
    # Spostiamo il file originale nella cartella archivio rinominandolo con il timestamp
    nome_archivio = f"Export_{datetime.now().strftime('%Y%m%d_%H%M%S')}.csv"
    percorso_destinazione = os.path.join(ARCHIVE_DIR, nome_archivio)

    shutil.move(CSV_FILE, percorso_destinazione)
    print(f"[{datetime.now()}] File originale archiviato con successo in: {percorso_destinazione}")

except Exception as e:
    print(f"[{datetime.now()}] ATTENZIONE: Mail inviata, ma archiviazione del file fallita. Errore: {e}")
\end{lstlisting}




%TO CHECK THE CONTENT

\chapter{Monitoraggio host non raggiungibili e mai raggiunti smartCity Edison}

Per garantire l'integrità del sistema di monitoraggio e per motivi di manutenzione,
è fondamentale non solo rilevare quando un dispositivo smette di rispondere, 
ma anche identificare eventuali discrepanze tra l'inventario .xlsx (apparati installati) e l'inventario logico presente su Zabbix. 
A tale scopo è stato implementato un sistema di analisi incrociata 
che confronta i dati estratti in tempo reale dalle API di Zabbix con un database censito in formato Excel.

\section{Descrizione dell'architettura}
Il sistema si pone l'obiettivo di generare un report su due categorie di criticità:
\begin{itemize}
    \item \textbf{Host non raggiungibili}: dispositivi presenti su Zabbix che hanno trigger attivi per il mancato ICMP ping.
    \item \textbf{Host mai raggiunti}: dispositivi presenti nel file di inventario Excel ma mai registrati nel monitoraggio Zabbix.
\end{itemize}

Il workflow è suddiviso in tre componenti principali operanti nell ambiente virtuale Python (venv):
\begin{enumerate}
    \item \textbf{Script di estrazione (\texttt{Python})}: analizza Zabbix e incrocia i dati con l'Excel.
    \item \textbf{Manager operativo (\texttt{Python})}: gestisce l'interazione con l'utente e l'invio delle email.
    \item \textbf{Wrapper di avvio (\texttt{Bash})}: semplifica l'esecuzione per utenti non tecnici garantendo il corretto caricamento dell'ambiente.
\end{enumerate}

\section{Script di estrazione e analisi}
Il cuore logico del sistema risiede nel file: \\
\inlinecode{/home/administrator/zabbixHostsScripts/notReachableHostsScript/}\\
\inlinecode{estraiHostNonRaggiungibili_e_maiRaggiunti.py}.

Lo script esegue le seguenti operazioni:
\begin{itemize}
    \item Si connette alle API di Zabbix tramite Token e recupera gli host appartenenti al gruppo principale \texttt{smartCity} e relativi sottogruppi.
    \item Carica il file di inventario situato nella directory superiore: \\ \inlinecode{/home/administrator/zabbixHostsScripts/HostInfosUbuntu.xlsx}.
    \item Esegue una funzione di normalizzazione degli indirizzi IP (\texttt{correggi\_ip\_excel}) per gestire formattazioni errate provenienti dall'Excel.
    \item Filtra i problemi attivi cercando specificamente la stringa: \\ \texttt{ICMP Ping: Unavailable by ICMP ping}.
    \item Genera un file CSV nella cartella \texttt{reports/} nominato con timestamp (es. \texttt{20240522\_103000\_report\_host\_non\_raggiungibili.csv}).
\end{itemize}

\section{Gestione operativa: run\_manager.py}
Per rendere il sistema interattivo, è stato implementato un middleware di gestione: \\
\inlinecode{/home/administrator/zabbixHostsScripts/notReachableHostsScript/run_manager.py}.

Questo script funge da interfaccia utente e permette di:
\begin{itemize}
    \item \textbf{Esclusione Dinamica}: Chiedere all'operatore quali comuni escludere dall'analisi (es. "MUSC", "BORG").
    \item \textbf{Identificazione Report}: Individuare automaticamente l'ultimo file CSV generato nella cartella \texttt{reports/}.
    \item \textbf{Distribuzione}: Chiedere gli indirizzi email a cui inviare il report finale.
\end{itemize}

\subsection{Configurazione e invio email}
L'invio avviene tramite un servizio python interno che non richiede autenticazione, configurato per comunicare sulla porta 25.

\begin{lstlisting}[language=python, caption=Configurazione mail in run\_manager.py]
# Configurazione per invio tramite server interno
SMTP_SERVER = "pmg.secoval.it"
SMTP_PORT = 25
MITTENTE_EMAIL = "zabbix@secoval.it"
\end{lstlisting}

\section{Automazione e accessibilità (bash wrapper)}
Per facilitare l'esecuzione da parte di personale amministrativo senza competenze tecniche di terminale è stato creato uno script Bash: \\
\inlinecode{/home/administrator/zabbixHostsScripts}\\
\inlinecode{/notReachableHostsScript/export_dei_punti_che_non_funzionano.sh}.

Questo script automatizza il cambio di directory e l'attivazione silenziosa del Virtual Environment, garantendo che i percorsi relativi (Excel e Reports) siano sempre risolti correttamente.

\begin{lstlisting}[language=bash, caption=Contenuto del wrapper avvia\_report.sh]
#!/bin/bash

# 1. Ci spostiamo nella cartella di lavoro corretta.
# Questo è vitale affinché lo script Python trovi l'Excel e la cartella reports.
cd /home/administrator/zabbixHostsScripts/notReachableHostsScript/ || {
    echo "Errore: Impossibile trovare la cartella dello script."
    exit 1
}

# 2. Avviamo il manager usando il Python del Virtual Environment
/home/administrator/zabbixHostsScripts/venv/bin/python3 run_manager.py
\end{lstlisting}

\subsection{Esecuzione globale}
Lo script è stato reso eseguibile tramite il comando \texttt{chmod +x} ed entrando in SSH nella VM ubuntu esso 
può essere richiamato da qualsiasi posizione nel sistema utilizzando il percorso assoluto:
\begin{lstlisting}[language=bash]
/home/administrator/zabbixHostsScripts/notReachableHostsScript/export_dei_punti_che_non_funzionano.sh
\end{lstlisting}

\section{Struttura delle directory}
La gerarchia dei file è organizzata come segue:

\begin{verbatim}
/home/administrator/zabbixHostsScripts/
├── cron_log.log
├── dailyRecordsScripts
│    ├── emailAndClearCsvProblemFile_mailer_log.log
│    ├── emailAndClearCsvProblemFile.py
│    └── zabbixTriggeredEventsArchive
│        ├── Export_20260429_150735.csv
│        ├── Export_20260430_083001.csv
│        ├── ...
├── hikvisionOscuramentoScripts
│    ├── config_telecamere_log.log
│    ├── configura_telecamere.sh
│    ├── estraiIPtelecamereHikvision.py
│    ├── hikvision_listener_log.log
│    ├── hikvision_listener_multiplehosts.py
│    ├── ipHik_log.log
│    ├── ip_hikvision.txt
│    ├── storico_log
│    │   ├── config_telecamere_log.log-2026-04-09.gz
│    │   ├── ...
│    │   ├── hikvision_listener_log.log-2026-04-09.gz
│    │   ├── ...
│    │   ├── ipHik_log.log-2026-04-09.gz
│    │   ├── ...
│    ├── test_ip_hikvision.txt
│    └── XMLs
│        ├── oscuramento_setup.xml
│        ├── triggers_oscuramento_setup.xml
│        ├── triggers_VideoMotionDetection_setup.xml
│        └── VideoMotionDetection_setup.xml
├── HostInfosUbuntu.xlsx
├── info_e_location_pali_ubuntu.xlsx
├── notReachableHostsScript
│    ├── estraiHostNonRaggiungibili_e_maiRaggiunti.py
│    ├── export_dei_punti_che_non_funzionano.sh
│    ├── reports
│    │   ├── 20260422_145058_report_host_non_raggiungibili.csv
│    │   ├── 20260422_150946_report_host_non_raggiungibili.csv
│    │   ├── ...
│    └── run_manager.py
├── storico_log
│    ├── cron_log.log-2026-03-09.gz
│    ├── cron_log.log-2026-03-15.gz
│    ├── ...
│    └── deprecated_cron_pali_info.log
├── venv
│    ├─ ...
├── zabbix_HostName_IDpalo_updater.py
└── zabbix_palo_location_and_info_updater.py
\end{verbatim}

\chapter{Analisi e calcolo dei KPI}

\section{Scopo dell'analisi storica}
Come illustrato nei capitoli precedenti, 
il sistema esporta quotidianamente i dati relativi al ciclo di vita dei \emph{Problem} 
generati da Zabbix all'interno di file CSV archiviati nella directory:
\inlinecode{/home/administrator/zabbixHostsScripts/dailyRecordsScripts}\\\inlinecode{/zabbixTriggeredEventsArchive}.

Sebbene i report via email offrano una fotografia giornaliera, 
è emersa la necessità di estrarre \textbf{Key Performance Indicators (KPI)} su scala mensile e annuale. 
L'obiettivo è fornire alla direzione tecnica metriche di alto livello, 
quali il tempo medio di risoluzione dei guasti (MTTR), 
la percentuale di host attualmente non funzionanti rispetto al totale installato e la distribuzione settimanale e oraria delle criticità.

Per soddisfare questa esigenza sono stati sviluppati due script Python, 
eseguiti in modo isolato tramite l'ambiente virtuale (\texttt{venv}).
Essi aggregano l'intero storico dei CSV e calcolano statistiche globali o filtrate per specifiche sottoreti.

\section{Script globale: analizeProblemsAndGetKPIs.py}
Il primo strumento è lo script globale \inlinecode{analizeProblemsAndGetKPIs.py}. Il suo compito è analizzare trasversalmente tutti i dispositivi dell'infrastruttura (TVCC, SW, AP, LTE).

\subsection{Logica di funzionamento}
Lo script utilizza la libreria \texttt{glob} per individuare tutti i file CSV presenti nell'archivio. Successivamente, elabora le righe estraendo informazioni vitali per i KPI:
\begin{itemize}
    \item \textbf{Tracciamento EventID:} Attraverso l'uso di insiemi (\texttt{set()}), lo script memorizza gli \texttt{EventID} con stato \texttt{PROBLEM} in un set di eventi aperti, e quelli con stato \texttt{RESOLVED} in un set di eventi chiusi. La differenza tra i due set (\inlinecode{eventi_aperti - eventi_chiusi}) restituisce matematicamente il numero esatto di dispositivi attualmente guasti.
    \item \textbf{Calcolo MTTR (Mean Time To Resolution):} Invece di eseguire complesse sottrazioni tra date, lo script sfrutta l'intelligenza di Zabbix leggendo direttamente la colonna \texttt{Duration}. Tramite un'espressione regolare (\texttt{re.finditer}), la stringa di durata (es. \texttt{4d 21h 54m 2s}) viene scomposta e convertita in secondi totali per permetterne la media aritmetica.
    \item \textbf{Categorizzazione delle Reti:} La funzione \texttt{categorizza\_rete(ip\_string)} divide statisticamente i problemi basandosi sul terzo e quarto ottetto dell'IP (es. \texttt{x.x.241.x} con quarto ottetto \texttt{0-127} viene classificato come TVCC).
\end{itemize}

\begin{lstlisting}[language=Python, caption=Estrazione dei secondi dalla colonna Duration tramite Regex]
def parse_duration_to_seconds(duration_str):
    if not duration_str:
        return 0
    secondi_totali = 0
    # Trova tutti i numeri seguiti da d, h, m, o s
    matches = re.finditer(r'(\d+)\s*([dhms])', duration_str.lower())
    
    for match in matches:
        valore = int(match.group(1))
        unita = match.group(2)
        if unita == 'd': secondi_totali += valore * 86400
        elif unita == 'h': secondi_totali += valore * 3600
        elif unita == 'm': secondi_totali += valore * 60
        elif unita == 's': secondi_totali += valore
            
    return secondi_totali
\end{lstlisting}

\subsection{Modalità di esecuzione}
Per garantire il corretto funzionamento delle librerie standard di Python e per non intaccare l'OS host, lo script deve essere eseguito utilizzando l'eseguibile Python situato nel \texttt{venv} di progetto.
Da riga di comando:

\begin{lstlisting}[language=bash, caption=Esecuzione dello script di analisi globale]
/home/administrator/zabbixHostsScripts/venv/bin/python3 /home/administrator/zabbixHostsScripts/dailyRecordsScripts/analizeProblemsAndGetKPIs.py
\end{lstlisting}

\section{Script Specifico: TVCCanalizeProblemsAndGetKPIs.py}
Per ottenere report più granulari, utili soprattutto a chi si occupa unicamente della videosorveglianza, è stato derivato un secondo script: \inlinecode{TVCCanalizeProblemsAndGetKPIs.py}.

\subsection{Differenze e Filtraggio}
La struttura dello script è pressoché identica al globale, con due differenze architetturali fondamentali:
\begin{enumerate}
    \item \textbf{Il "Grandissimo IF" (Filtro IP):} Durante il ciclo di lettura del CSV, l'IP viene immediatamente categorizzato. Se il risultato è diverso da \texttt{TVCC}, lo script esegue un \texttt{continue}, saltando l'elaborazione della riga. Questo garantisce che ogni singola metrica (MTTR, fasce orarie, top host) sia inquinata \textbf{esclusivamente} da telecamere.
    \item \textbf{Costante Hardcoded Mirata:} La variabile \inlinecode{NUMERO_HOST_TVCC_ZABBIX = 607} sostituisce il totale globale. In questo modo, l'incidenza percentuale dei problemi (es. il 3.4\% della rete è guasto) viene calcolata solo in base al numero reale di telecamere installate e non sull'intero parco macchine.
\end{enumerate}

\begin{lstlisting}[language=Python, caption=Il blocco logico di filtraggio della sottorete TVCC]
for row in reader:
    ip = row.get('IP', '').strip()

    # Se l'IP non appartiene alla rete TVCC, salta la riga
    if categorizza_rete(ip) != "TVCC":
        continue

    # Da qui in poi, lo script elabora SOLO dati TVCC
    subnet_counter["TVCC"] += 1
    # ... segue estrazione di EventID e Status ...
\end{lstlisting}

\subsection{Modalità di esecuzione}
Analogamente allo script globale, l'esecuzione avviene tramite l'ambiente virtualizzato:

\begin{lstlisting}[language=bash, caption=Esecuzione dello script di analisi specifica TVCC]
/home/administrator/zabbixHostsScripts/venv/bin/python3 /home/administrator/zabbixHostsScripts/dailyRecordsScripts/TVCCanalizeProblemsAndGetKPIs.py
\end{lstlisting}

\section{Descrizione dell'Output}
Entrambi gli script, una volta eseguiti, generano uno \textit{standard output} formattato in puro testo ASCII (ideale per la lettura rapida su terminale), diviso nelle seguenti sezioni analitiche:

\begin{description}
    \item[Metriche Chiave (KPI):] Mostra il numero assoluto di problemi "aperti" e la loro incidenza percentuale rispetto alla costante \texttt{NUMERO\_HOST\_ZABBIX}. Mostra inoltre l'MTTR convertito in giorni, ore, minuti e secondi.
    \item[Analisi delle Sottoreti:] Conferma la distribuzione degli eventi sulle varie VLAN (o conferma il 100\% su TVCC nel caso dello script specifico).
    \item[Classifica TOP Host:] Genera una \textit{top ten} degli apparati più problematici (utile per individuare dispositivi in \textit{flapping} costante).
    \item[Analisi Temporale:] Due blocchi separati per individuare i \textit{pattern} temporali: un istogramma ASCII diviso in fasce orarie e un conteggio organizzato per singola settimana (ISO week).
\end{description}

\chapter{Integrazione ticketing Interacta}

\section{Scopo e architettura dell'integrazione}
Per garantire l'integrazione con il sistema già presente e il tracciamento della gestione e degli interventi fisici 
sugli apparati della Smart City, 
si è resa necessaria un'integrazione tra il sistema di monitoraggio Zabbix 
e il portale (ITSM, IT Service Management) \textbf{Interacta}.

L'obiettivo dell'integrazione è automatizzare l'apertura di un ticket di assistenza
(con tutti i dettagli dell'host, le coordinate e il tipo di guasto) 
nel momento in cui un trigger di Zabbix entra in stato \texttt{PROBLEM}, 
e chiuderlo in automatico quando lo stato torna a \texttt{RESOLVED}.

L'architettura si basa su tre pilastri:
\begin{itemize}
    \item Un'autenticazione \textit{Server-to-Server} verso le API di Interacta tramite Service Account 
    (Service Account Authentication è un metodo che consente ad applicazioni, script o macchine virtuali
    di verificare in modo sicuro la propria identità e interagire con API o servizi cloud senza intervento umano) e
    firma crittografica asimmetrica (RSA).
    \item Uno script Python (\texttt{interactaCreateTicket.py}) operante in 
    \\\texttt{/usr/lib/zabbix/alertscripts/}, incaricato di elaborare i dati e comunicare con le API.
    \item Un Media Type di tipo "Script" configurato su Zabbix per inoltrare i parametri 
    formattati allo script in tempo reale.
\end{itemize}

\section{L'Autenticazione Server-to-Server (Il JWT manuale)}
L'aspetto tecnicamente più complesso dell'integrazione ha riguardato l'ottenimento 
del \textit{Bearer Token} per interrogare le API di Interacta.

La documentazione fornita per il Service Account richiedeva la generazione di una \textit{JWT Assertion}(Json Web Token) 
firmata con algoritmo \textbf{RS512} utilizzando una chiave privata fornita in formato JSON (\texttt{service-account-16-prod.json}). 
Tuttavia, è emersa un'incongruenza critica tra lo standard mondiale dell'ingegneria dei JWT (RFC 7519) 
e l'implementazione del backend di Interacta: il backend richiede esplicitamente che il campo
\texttt{kid} (Key ID) nell'Header del token sia un \textbf{numero intero} 
(es. \texttt{16}), mentre lo standard (e le librerie Python rigorose come \texttt{PyJWT}) esigono che sia una \textbf{stringa}.

Per evitare errori dovuti al rifiuto di \texttt{PyJWT} di inserire un intero nell'Header, 
si è optato per la \textbf{costruzione e firma manuale del JWT} 
utilizzando direttamente la libreria di basso livello \texttt{cryptography}.

\begin{lstlisting}[language=Python, caption={Costruzione manuale del JWT per forzare il campo 'kid' come intero}]
header = {
    "kid": key_id, # MANTENUTO COME NUMERO INTERO PER IL BACKEND INTERACTA
    "alg": "RS512"
}
payload = {
    "jti": str(uuid.uuid4()),
    "aud": "injenia/portal-authenticator",
    "iss": str(client_id),
    "iat": now - 60, # Offset per prevenire errori di time-drift
    "exp": now + 600 # Scadenza rigida (massimo 10 minuti)
}

# Codifica in Base64 URL-safe ed esecuzione manuale della firma RSA512
b64_header = b64url_encode(json.dumps(header).encode('utf-8'))
b64_payload = b64url_encode(json.dumps(payload).encode('utf-8'))
signing_input = f"{b64_header}.{b64_payload}".encode('utf-8')

private_key_obj = load_pem_private_key(private_key_pem.encode('utf-8'), password=None)
signature = private_key_obj.sign(signing_input, padding.PKCS1v15(), hashes.SHA512())

signed_jwt = f"{b64_header}.{b64_payload}.{b64url_encode(signature)}"
\end{lstlisting}

\section{Creazione del ticket e formattazione della payload}
Una volta ottenuto l'access token, lo script formatta i dati passati da Zabbix per inviarli all'endpoint specifico 
\\(\texttt{/extensions/SmartCity/richiesta}).

Durante lo sviluppo, è stato cruciale rispettare rigorosamente i tipi di dato richiesti dal database
di Interacta per evitare errori \texttt{500 Internal Server Error}:
\begin{itemize}
    \item \textbf{Coordinate}: devono essere passate come stringa es. \texttt{"45.0703 7.6869"}.
    \item \textbf{Comune}: il backend si aspetta obbligatoriamente un Array di numeri interi, normalmente di dim=1, 
    non una stringa testuale (es. \texttt{[1]}). 
    Sono definiti da una specifica funzione in base al loro host\_name che contiene il codice a 4 
    caratteri per l'identificazione del comune 
    (tranne Castenedolo a cui è associato un codice a5 caratteri: "CASTE").
    \item \textbf{PROBLEM ID}: l'\texttt{\{EVENT.ID\}} di Zabbix deve essere castato in formato intero.
\end{itemize}

\begin{lstlisting}[language=Python, caption={Costruzione del payload per l'apertura del ticket}]
payload = {
    "fields": {
        "progetto": ticket_data["progetto"],
        "nome_host": ticket_data["nome_host"],
        "IP": ticket_data["IP"],
        "IDpalo": ticket_data["ID_palo"],
        "tipo_problema": ticket_data["tipo_problema"],
        "coordinate": f"{ticket_data["coordinate_lat"]} {ticket_data["coordinate_lon"]}",
        "via": ticket_data["via"],
        "PROBLEM ID": int(ticket_data["problem_ID"]),
        "comune": findComuneID(ticket_data["nome_host"])
    }
}
\end{lstlisting}

\textit{Nota importante}: la funzione \textbf{findComuneID} prende come input il nome\_host e restituisce un identificativo del comune,
attraverso il riferimento esterno di un post nella community Enti di Interacta.
Essa crea un dizionario a partire da un file CSV \inlinecode{/usr/lib/zabbix/alertscripts/post_export-202605131241_anagrafiche_enti.csv}.
Da qui viene estratta la colonna con gli "ID" e "Sigla Zabbix".
Vengono estratti i primi caratteri dall'host\_name prima del "-" e estratto l'ID relativo a quel codice.
Per ogni host\_name che non è riconosciuto dalla funzione il programma si interrompe loggando l'accaduto.
Questo permette di escudere tutti gli host come firewall e gateway.
\textbf{Ogni nuovo host fuori da progetto Edison necessita di essere inserito in zabbix
con il codice del comune come prefisso seguito da "-" (es.TBRE-TVCC-08 ) e per ogni nuovo comune/ente 
deve essere aggiornato il file csv; per il corretto funzionamento di questa feature}. 


\section{Risoluzione del ticket e formattazione della payload}
L'operazione di chiusura (tramite la verifica del flag \texttt{status == "RESOLVED"}) 
è predisposta nello script, in base agli argomenti ricevuti in input ed è una chiamata più semplice rispetto all'apertura.

\begin{lstlisting}[language=Python, caption={Costruzione del payload per la chiusura del ticket}]
    elif status == "RESOLVED":
        url = f"{API_BASE}/extensions/SmartCity/richiesta/chiusura" # UPDATE LATER: Final Create API endpoint
        payload = {
            "clientUid": int(ticket_data["problem_ID"])
        }
        with open(DEBUG_AND_LOG_FILE, "a") as f:
            f.write(f"ticket chiusura PAYLOAD created\n")

        r = requests.post(url, headers=headers, json=payload, allow_redirects=True)
        if r.status_code == 500:
            raise HTTPError("Returned status code: 500. Probably tried to close an already closed ticket", response=r)
        r.raise_for_status()
}
\end{lstlisting}

\section{Configurazione lato Zabbix}
Affinché Zabbix possa lanciare lo script al momento giusto, 
sono stati configurati diversi elementi sulla piattaforma. 
Se in futuro l'integrazione dovesse smettere di funzionare, questa è la \textit{checklist} primaria da verificare.

\subsection{Media type (Interacta ticket script)}
In \textbf{Alerts $\rightarrow$ Media types} è stato creato un Media di tipo \textit{Script} 
puntato al file \texttt{/usr/lib/zabbix/alertscripts/interactaCreateTicket.py}.
Nella sezione \textit{Script parameters} sono mappate in modo rigido le macro di Zabbix in flag
utilizzabili dalla libreria \texttt{argparse} di Python:
\begin{itemize}
    \item \texttt{--status} $\rightarrow$ \texttt{\{EVENT.STATUS\}}
    \item \texttt{--progetto} $\rightarrow$ \texttt{\{TRIGGER.HOSTGROUP.NAME\}}
    \item \texttt{--nome\_host} $\rightarrow$ \texttt{\{HOST.NAME\}}
    \item \texttt{--IP} $\rightarrow$ \texttt{\{HOST.IP\}}
    \item \texttt{--ID\_palo} $\rightarrow$ \texttt{\{INVENTORY.LOCATION\}} (o campo equivalente)
    \item \texttt{--tipo\_problema} $\rightarrow$ \texttt{\{EVENT.NAME\}}
    \item \texttt{--coordinate\_lat} $\rightarrow$ \texttt{\{INVENTORY.LOCATION.LAT\}}
    \item \texttt{--coordinate\_lon} $\rightarrow$ \texttt{\{INVENTORY.LOCATION.LON\}}
    \item \texttt{--via} $\rightarrow$ \texttt{\{INVENTORY.SITE.ADDRESS.A\}}
    \item \texttt{--problem\_ID} $\rightarrow$ \texttt{\{EVENT.ID\}}
\end{itemize}

% \textbf{Attenzione - Il vincolo dei Message Templates:} Zabbix si rifiuta di eseguire Media Type personalizzati (restituendo l'errore \textit{"No message defined for media type"}) 
% se non vengono definiti dei template testuali fittizi. 
% È stato quindi necessario configurare nella tab \textit{Message templates} dei template generici per \texttt{Problem} e \texttt{Problem recovery}, anche se lo script Python li ignora completamente.

\subsection{Trigger actions e permessi utente}
L'esecuzione dello script è agganciata a diversi \textbf{Trigger} già
presenti nel sistema per lanciare alarmi riguardo il ping e l'oscuramento.

% \textbf{Nota sui permessi:} se l'invio è assegnato a un utente non-SuperAdmin (es. un utente di servizio \texttt{api\_interacta}), è \textit{tassativo} 
% che il gruppo a cui appartiene l'utente abbia i privilegi di \textbf{Lettura (Read)} 
% sugli Host Group monitorati (in \textit{Users $\rightarrow$ User groups $\rightarrow$ Host permissions}). In caso contrario, l'azione fallirà silenziosamente poiché l'utente non "vede" l'host guasto.

\section{Gestione dei log e manutenzione (Logrotate)}
Lo script genera output di debug essenziali nel file 
\\\texttt{/usr/lib/zabbix/alertscripts/interactaCreateOrSolveTicket\_log.log}. 
A causa del contesto di esecuzione in background di Zabbix, 
lo script Python utilizza \textbf{percorsi assoluti} 
(es. \texttt{/usr/lib/zabbix/alertscripts/...}) per evitare errori di \texttt{PermissionDenied} sulla radice del filesystem.

Per evitare l'esaurimento dello spazio su disco, 
è stata implementata una configurazione di \textbf{Logrotate} dedicata, situata in \texttt{/etc/logrotate.d/zabbix-interacta-script}:

\begin{lstlisting}[language=bash, caption={Configurazione logrotate per i file di log dell'integrazione}]
/usr/lib/zabbix/alertscripts/interactaCreateOrSolveTicket_log.log {
    su zabbix zabbix
    weekly
    maxsize 30M
    rotate 12
    dateext
    dateformat -%Y-%m-%d
    olddir /usr/lib/zabbix/alertscripts/storico_log
    compress
    delaycompress
    missingok
    notifempty
    create 0600 zabbix zabbix
}
\end{lstlisting}
\textit{Nota architetturale:} Affinché la direttiva \texttt{su zabbix zabbix} 
potesse spostare i file ruotati all'interno della \texttt{olddir}, 
è stato assegnato il gruppo \texttt{zabbix} alla cartella padre \texttt{alertscripts} 
e sono stati concessi i permessi di scrittura al gruppo (\texttt{chmod g+w}).

\chapter{Quick reference: funzionalità e sistemi}
Il presente capitolo è un riassunto strutturato delle funzionalità attive, dei sistemi, degli script e dei relativi log.

\section{Automazione dell'inventario (metadati)}
Per arricchire i dispositivi monitorati con informazioni geografiche e tecniche, sono stati sviluppati script per l'aggiornamento automatico dei campi di inventory.
\begin{itemize}
    \item \textbf{Funzionalità}: estrazione dei dati da fogli Excel (nome host, ID del palo, coordinate geografiche, indirizzo e tecnologia)
     per arricchire automaticamente i campi di inventario in Zabbix mano a mano che i dispositivi vengono scoperti.
    \item \textbf{Script coinvolti}:
    \begin{itemize}
        \item \texttt{zabbix\_HostName\_IDpalo\_updater.py}.
        \item \texttt{zabbix\_palo\_location\_and\_info\_updater.py}.
    \end{itemize}
    \item \textbf{Sorgenti}: \texttt{Host\_InfosUbuntu.xlsx} e \\\texttt{info\_e\_location\_pali\_ubuntu.xlsx}.
    \item \textbf{Log}: \texttt{/home/administrator/zabbixHostsScripts/cron\_log.log} \\(sottoposto a rotazione tramite logrotate configurato in \\\texttt{/etc/logrotate.d/zabbix\_scripts}).
\end{itemize}

\section{Acquisizione dati hardware tramite template}
Riconoscimento automatico e acquisizione di informazioni diagnostiche e di modello interrogando i dispositivi in base al loro brand o tipologia.
\begin{itemize}
    \item \textbf{Hikvision}: utilizzo di chiamate HTTP ISAPI per estrarre modello e firmware.
    \item \textbf{Standard SNMP}: query SNMP per apparati di rete come Switch e Access Point.
    \item \textbf{Script coinvolti}: \texttt{createDiscoveryRulesForHikvisions.py}, \\\texttt{apCreateDiscoveryRules.py} e \texttt{swCreateDiscoveryRules.py}.
\end{itemize}

\section{Sicurezza e rilevamento manomissioni}
È stato implementato un sistema per rilevare eventi di \textit{tamper detection} (oscuramento obiettivo) e \textit{motion detection}.
\begin{itemize}
    \item \textbf{Configurazione XML}: setup remoto delle telecamere tramite file XML via HTTP PUT.
    \item \textbf{Listener Asincrono}: middleware che ascolta lo streaming degli eventi dalle telecamere e invia i dati a Zabbix tramite \texttt{zabbix\_sender}.
    \item \textbf{Script coinvolti}:
    \begin{itemize}
        \item \texttt{estraiIPtelecamereHikvision.py} (estrae da Zabbix gli IP delle telecamere)
        \item \texttt{configura\_telecamere.sh} (carica i file XML di configurazione tramite chiamata PUT asincrona)
        \item \texttt{hikvision\_listener\_multiplehosts.py} (middleware avviato come servizio demone che ascolta il flusso dati e invia gli eventi usando zabbix\_sender)
    \end{itemize}
    \item \textbf{Log}: \texttt{ipHik\_log.log}, \texttt{config\_telecamere\_log.log} \\e \texttt{hikvision\_listener\_log.log}.
\end{itemize}

\section{Reporting e analisi giornaliera}
Intercettazione degli allarmi attivi/risolti, esportazione del ciclo di vita del problema su un CSV locale e produzione di report giornalieri.
\begin{itemize}
    \item \textbf{Esportazione CSV}: gli eventi vengono salvati in locale tramite un Custom AlertScript.
    \item \textbf{Analisi Python}: elaborazione dei dati per identificare host instabili (\textit{flapping}) o allarmi di sicurezza critici.
    \item \textbf{Report mensili}: invio di report mensili con KPIs.
    \item \textbf{Script coinvolti}:
    \begin{itemize}
        \item \texttt{ScriptCustomExportProblemsToCSV.sh}.
        \item \texttt{emailAndClearCsvProblemFile.py}.
        \item \texttt{monthlyReport.py}
    \end{itemize}
    \item \textbf{Log}: \texttt{zabbix\_csv\_mailer.log}.
\end{itemize}

\section{Monitoraggio degli host "non raggiungibili" o "mai raggiunti"}
Il sistema confronta l'inventario fisico su Excel con quello logico su Zabbix per identificare host non raggiungibili o mai registrati.
\begin{itemize}
    \item \textbf{Analisi Incrociata}: tra gli apparati fisicamente installati (segnati su Excel) e lo stato del monitoraggio logico in Zabbix,
     generando un report sulle discrepanze e sui dispositivi offline.
    \item \textbf{Script coinvolti}:
    \begin{itemize}
        \item \texttt{estraiHostNonRaggiungibili\_e\_maiRaggiunti.py}.
        \item \texttt{run\_manager.py} (interfaccia utente interattiva).
        \item \texttt{export\_dei\_punti\_che\_non\_funzionano.sh} \\(wrapper per l'esecuzione).
    \end{itemize}
\end{itemize}

\section{Analisi e calcolo dei KPI}
Elaborazione dello storico degli eventi Zabbix per estrarre metriche di alto livello, su scala mensile e annuale.
\begin{itemize}
    \item \textbf{Funzionalità}: aggregazione dei dati esportati nei CSV per il tracciamento dei dispositivi attualmente guasti, 
    calcolo del \textit{Mean Time To Resolution} 
    (tramite regex sulla colonna \texttt{Duration}) e categorizzazione temporale e per sottorete.
    \item \textbf{Script coinvolti}:
    \begin{itemize}
        \item \texttt{analizeProblemsAndGetKPIs.py} (analisi globale di tutti gli apparati dell'infrastruttura).
        \item \texttt{TVCCanalizeProblemsAndGetKPIs.py} (analisi filtrata esclusivamente per le telecamere,
         con costanti hardware).
    \end{itemize}
    \item \textbf{Sorgenti}: archivio dei file CSV situato in 
    \\\texttt{/home/administrator/zabbixHostsScripts/dailyRecordsScripts/}\\\texttt{zabbixTriggeredEventsArchive}.
    \item \textbf{Esecuzione}: script eseguiti isolati tramite ambiente virtuale Python (\texttt{venv}).
    \item \textbf{Output}: report testuale ASCII generato in \textit{standard output} 
    (mostra metriche chiave, classifica TOP host e pattern temporali).
\end{itemize}

\section{Integrazione ticketing interacta}
Automazione dell'apertura e chiusura dei ticket di assistenza sul portale Interacta 
in sincronia con lo stato (\texttt{PROBLEM} o \texttt{RESOLVED}) dei trigger di Zabbix.
\begin{itemize}
    \item \textbf{Funzionalità}: autenticazione Server-to-Server tramite costruzione e firma manuale di un token JWT 
    (per requisiti specifici del backend Interacta). Lo script intercetta le macro Zabbix 
    (coordinate, comune, ID palo) e invia il payload formattato alle API del portale.
    \item \textbf{Script coinvolti}:
    \begin{itemize}
        \item \texttt{interactaCreateTicket.py} (posizionato e operante nella directory \\\texttt{/usr/lib/zabbix/alertscripts/}).
    \end{itemize}
    \item \textbf{Configurazione Zabbix}: l'integrazione sfrutta un \textit{Media type} tipo "Script" 
    che mappa le macro di Zabbix sugli argomenti dello script, richiamato da apposite \textit{Trigger actions}.
    \item \textbf{Credenziali}: chiave privata \texttt{service-account-16-prod.json} usata per la firma RSA512.
    \item \textbf{Log}: \texttt{/usr/lib/zabbix/alertscripts/interactaCreateOrSolveTicket}\\\texttt{\_log.log} 
    (sottoposto a rotazione tramite logrotate dedicato, configurato in \\\texttt{/etc/logrotate.d/zabbix-interacta-script}).
\end{itemize}

\end{document}



